% ================== ЛАБОРАТОРНАЯ РАБОТА №4 ==================
\documentclass[14pt,a4paper]{extarticle}

\usepackage[utf8]{inputenc}
\usepackage[russian]{babel}
\usepackage{geometry}
\geometry{left=30mm,right=15mm,top=20mm,bottom=20mm}

\usepackage{indentfirst}
\setlength{\parindent}{1.25cm}

\usepackage{amsmath,amssymb}
\usepackage{graphicx}
\usepackage{float}
\usepackage{xcolor}
\usepackage{hyperref}
\usepackage{caption}
\usepackage{setspace}
\setstretch{1.0}

\usepackage{titlesec}
\titleformat{\section}{\normalfont\normalsize\bfseries}{\thesection}{1em}{\MakeUppercase}
\titleformat{\subsection}{\normalfont\normalsize\bfseries}{\thesubsection}{1em}{}
\titlespacing{\section}{\parindent}{3.5ex plus 1ex minus .2ex}{2.3ex plus .2ex}
\titlespacing{\subsection}{\parindent}{3.25ex plus 1ex minus .2ex}{1.5ex plus .2ex}

\usepackage{fancyhdr}
\pagestyle{fancy}
\fancyhf{}
\fancyfoot[R]{\thepage}
\renewcommand{\headrulewidth}{0pt}
\fancypagestyle{plain}{
  \fancyhf{}
  \fancyfoot[R]{\thepage}
  \renewcommand{\headrulewidth}{0pt}
}

\renewcommand{\contentsname}{\centering СОДЕРЖАНИЕ}

\begin{document}

\thispagestyle{empty}

\begin{center}
МИНИСТЕРСТВО ОБРАЗОВАНИЯ РЕСПУБЛИКИ БЕЛАРУСЬ\\[0.5em]
Учреждение образования\\[0.5em]
\textbf{<<Белорусский государственный университет информатики и радиоэлектроники>>}\\[2em]

Факультет компьютерных систем и сетей\\[0.5em]
Кафедра программного обеспечения информационных технологий
\end{center}

\vspace{3em}

\begin{center}
\textbf{ОТЧЁТ}\\[0.5em]
по лабораторной работе №4\\[0.5em]
по дисциплине\\[0.5em]
\textbf{<<Генетические и эволюционные вычисления>>}

\textbf{Тема: <<Кластеризация с помощью нейронной сети Кохонена>>}
\end{center}

\vspace{4em}

\hfill\begin{minipage}{0.55\textwidth}
Выполнил:\\
магистрант 1 курса магистратуры\\
факультета компьютерных систем и сетей\\
группы 556301\\[0.3em]
\textbf{Лебедевич Артём Владимирович}\\[1.2em]
Преподаватель:\\
доцент кафедры ПОИТ, канд. техн. наук\\
\textbf{Нестеренков Сергей Николаевич}
\end{minipage}

\vfill

\begin{center}
Минск, 2026~г.
\end{center}

\newpage

\tableofcontents

\newpage

\section{Цель работы}

Изучить принципы построения и обучения самоорганизующихся карт Кохонена (SOM) для задачи кластеризации. Реализовать сеть Кохонена для разделения множества точек на декартовой плоскости на заданное число кластеров.

\section{Теоретические сведения}

\subsection{Самоорганизующаяся карта Кохонена}

Самоорганизующаяся карта (Self-Organizing Map, SOM)~--- это тип нейронной сети, предложенный Т.~Кохоненом в 1982~г., реализующий обучение без учителя. Сеть состоит из $M$ нейронов, каждый из которых характеризуется вектором весов $\mathbf{w}_j \in \mathbb{R}^d$, $j = 1, \ldots, M$.

\subsection{Алгоритм обучения}

\begin{enumerate}
\item \textbf{Поиск нейрона-победителя (BMU)}. Для каждого входного вектора $\mathbf{x}$ находится нейрон с ближайшим вектором весов:
\begin{equation}
c = \arg\min_{j} \|\mathbf{x} - \mathbf{w}_j\|.
\end{equation}

\item \textbf{Обновление весов}. Веса BMU и его соседей обновляются:
\begin{equation}
\mathbf{w}_j^{(t+1)} = \mathbf{w}_j^{(t)} + \alpha(t) \cdot h_{cj}(t) \cdot (\mathbf{x} - \mathbf{w}_j^{(t)}),
\end{equation}

где $\alpha(t)$~--- скорость обучения, $h_{cj}(t)$~--- функция соседства.

\item \textbf{Функция соседства} (гауссова):
\begin{equation}
h_{cj}(t) = \exp\left(-\frac{\|r_c - r_j\|^2}{2\sigma(t)^2}\right),
\end{equation}

где $r_c, r_j$~--- позиции нейронов, $\sigma(t)$~--- радиус соседства.

\item \textbf{Затухание параметров}:
\begin{equation}
\alpha(t) = \alpha_0 \cdot \exp(-t/\tau), \qquad
\sigma(t) = \sigma_0 \cdot \exp(-t/\tau).
\end{equation}
\end{enumerate}

\subsection{Кластеризация на основе SOM}

После обучения каждому входному вектору присваивается метка нейрона-победителя. Таким образом, $M$ нейронов определяют $M$ кластеров. Качество кластеризации оценивается с помощью коэффициента силуэта:
\begin{equation}
s(i) = \frac{b(i) - a(i)}{\max(a(i), b(i))},
\end{equation}

где $a(i)$~--- среднее расстояние до точек своего кластера, $b(i)$~--- минимальное среднее расстояние до точек ближайшего чужого кластера.

\section{Ход работы}

\subsection{Генерация данных}

Сгенерировано 200 точек на плоскости, организованных в 4 кластера с центрами $(2, 2)$, $(-2, 2)$, $(2, -2)$, $(-2, -2)$ и стандартным отклонением $\sigma = 0.5$.

\subsection{Реализация сети Кохонена}

Реализована сеть Кохонена:
\begin{itemize}
\item Инициализация весов случайным образом в диапазоне данных.
\item Поиск BMU по евклидову расстоянию.
\item Гауссова функция соседства с экспоненциальным затуханием.
\item Сохранение снимков позиций нейронов для визуализации обучения.
\end{itemize}

\subsection{Обучение}

Сеть обучена с параметрами: $M = 4$ нейрона, 200 эпох, $\alpha_0 = 0.5$, $\sigma_0 = 2.0$. Визуализация показывает, что нейроны перемещаются от начальных случайных позиций к центрам кластеров данных.

\subsection{Анализ параметров}

Исследовано влияние:
\begin{itemize}
\item \textbf{Числа нейронов} ($M = 3, 4, 5, 6$): при $M = 4$ достигается наилучший силуэтный коэффициент, что соответствует истинному числу кластеров. При $M < 4$ кластеры объединяются, при $M > 4$~--- происходит избыточное дробление.
\item \textbf{Скорости обучения} ($\alpha_0 = 0.1, 0.3, 0.5, 0.8$): слишком низкая скорость замедляет сходимость, слишком высокая может вызвать нестабильность.
\item \textbf{Числа эпох} (50, 100, 200, 500): сходимость достигается к 200 эпохам, дальнейшее обучение не улучшает результат.
\end{itemize}

\section{Результаты}

\begin{enumerate}
\item Сеть Кохонена успешно кластеризует 2D-данные, выявляя структуру из 4 кластеров.
\item Коэффициент силуэта при $M = 4$ составляет $\approx 0.85$, что свидетельствует о чётко выраженной кластерной структуре.
\item Оптимальное число нейронов совпадает с истинным числом кластеров.
\item Экспоненциальное затухание параметров обучения обеспечивает переход от глобальной к локальной настройке.
\end{enumerate}

\section{Выводы}

В ходе лабораторной работы реализована самоорганизующаяся карта Кохонена для задачи кластеризации. Экспериментально показано, что SOM способна выявлять кластерную структуру данных методом обучения без учителя. Качество кластеризации зависит от выбора числа нейронов, скорости обучения и количества эпох. Гауссова функция соседства с экспоненциальным затуханием обеспечивает плавное обучение и устойчивую сходимость.

\end{document}
